%%=============================================================================
%% Voorwoord
%%=============================================================================

\chapter*{\IfLanguageName{dutch}{Woord vooraf}{Preface}}%
\label{ch:voorwoord}

%% TODO:
%% Het voorwoord is het enige deel van de bachelorproef waar je vanuit je
%% eigen standpunt (``ik-vorm'') mag schrijven. Je kan hier bv. motiveren
%% waarom jij het onderwerp wil bespreken.
%% Vergeet ook niet te bedanken wie je geholpen/gesteund/... heeft

In deze bachelorproef heb ik ervoor gekozen om het onderwerp operationele veerkracht binnen mainframe-omgevingen te onderzoeken, met bijzondere aandacht voor replicatiearchitecturen en hun impact op dataconsistentie en beschikbaarheid.
De aanleiding voor dit onderzoek ligt in de toenemende relevantie van operationele veerkracht in een sterk gedigitaliseerde en geopolitiek instabiele wereld. 
Financiële instellingen spelen hierin een cruciale rol, aangezien zij instaan voor het vertrouwen en de stabiliteit van het economische systeem. 
Storingen of cyberincidenten binnen deze sector kunnen immers verstrekkende gevolgen hebben voor zowel organisaties als eindgebruikers.
Daarnaast is cyberveiligheid binnen financiële omgevingen een steeds belangrijker aandachtspunt geworden. 
Het garanderen van continue beschikbaarheid, correcte transactieverwerking en minimale dataverliesmomenten is essentieel om de integriteit van financiële systemen te waarborgen.
Tot slot wil ik mijn co-promotor Jan Cannaerts bedanken voor de begeleiding, feedback en ondersteuning tijdens het uitwerken van deze bachelorproef.